\begin{textsample}{NEG Dim 1 – 2013\_c4 – Score -5.00 – 2013\_c4/enade\_2013\_serv\_soc\_b0ca85ce74.txt}  \label{ex:f1_neg_056}
CAPÍTULO 4 PERCEPÇÃO DA PROVA As análises feitas neste capítulo tratam das percepções dos concluintes da Área de Serviço Social sobre a prova aplicada no ENADE/2013.
Estas percepções foram mensuradas por meio de nove questões que avaliaram desde o grau de dificuldade da prova até o tempo gasto para concluí -la.
As percepções sobre a prova foram relacionadas com o desempenho dos estudantes e com a Grande Região de funcionamento do curso.
O questionário de percepção da prova encontra -se ao final do Anexo IV, que traz a reprodução da prova.
O desempenho dos estudantes foi classificado em quatro quartos.
Para tanto, esse desempenho foi ordenado de forma ascendente.
O percentil 25, P25, também conhecido como primeiro quartil, é a nota de desempenho que deixa um quarto ( 25 \% ) dos valores observados abaixo e três quartos acima.
A Figura 4.1 apresenta uma ilustração deste conceito.
O quarto inferior de desempenho é composto pelas notas abaixo do primeiro quartil.
Já o percentil 75, P75, também conhecido como terceiro quartil, é o valor para o qual há três quartos ( 75 \% ) dos dados abaixo e um quarto acima dele.
O quarto superior de desempenho é composto pelas notas iguais ou acima do terceiro quartil.
O percentil 50, P50, também conhecido como mediana, é o valor que divide as notas em dois conjuntos de igual tamanho.
O segundo quarto inclui valores entre o primeiro quartil ( P25 ) e a mediana.
O terceiro quarto contém os valores entre a mediana ( P50 ) e o terceiro quartil ( P75 ).
Vale ressaltar que percentis, quartis e medianas são pontos que não obrigatoriamente pertencem ao conjunto original de dados, ao passo que os quartos são subconjuntos dos dados originais.
A seguir, serão apresentados gráficos com resultados selecionados, relativos às nove questões avaliadas por grupos de estudantes.
Os gráficos apresentam nas barras o percentual de alunos que assinalaram uma das opções ou a soma das porcentagens daqueles que assinalaram duas ( ou três ) delas.
Por exemplo, para as questões 1 e 2, os gráficos apresentam a porcentagem total de participantes que assinalaram as opções ( D ) difícil e ( E ) muito difícil.
Em \textbf{cada} barra foram assinalados também os extremos do intervalo de confiança de 95 \% como linhas verticais unidas por uma linha horizontal na forma da letra H maiúscula, semelhantemente aos gráficos do Capítulo 3.
As Tabelas no Anexo II apresentam os valores absolutos e a distribuição percentual13 das alternativas válidas das nove questões, segundo o mesmo recorte de desempenho dos alunos e Grande Região de funcionamento do curso. 4.1 GRAU DE DIFICULDADE DA PROVA 4.1.1 Componente de Formação Geral Ao avaliarem “ Qual o grau de dificuldade desta prova na Parte de Formação Geral? ” ( Questão 1 ), 31,4 \% do grupo de inscritos e presentes optaram pelas alternativas difícil ou muito difícil.
Entretanto, para mais da metade dos estudantes ( 60,2 \% ), o Componente de Formação Geral da prova foi considerado com grau de dificuldade médio ( Gráfico 4.1, Gráfico 4.2 e, no Anexo II, a Tabela II.1 ).
O percentual de estudantes que consideraram a prova como difícil ou muito difícil foi maior na região Sul, onde a proporção foi de 35,5 \%, enquanto a de menor incidência foi a Nordeste, com 21,3 \%.
No Gráfico 4.1 é possível observar que esta diferença é estatisticamente significativa.
Esta proporção na região Nordeste se mostrou estatisticamente diferente de todas as demais regiões.
Nas Grandes Regiões, a proporção de presentes à prova que consideraram o Componente de Formação Geral como sendo de grau de dificuldade médio \textbf{esteve} entre 58,1 \% ( região Sul ) e 66,8 \% ( região Nordeste ), situação inversa dos que consideraram a prova como difícil ou muito difícil. 13 Cumpre \textbf{lembrar} uma das convenções para tabelas numéricas ( pág. iii ) sobre a possibilidade da soma das partes não resultar em 100 \% por questões de arredondamento.
O percentual de alunos que consideraram a prova difícil ou muito difícil foi decrescente em relação ao aumento de desempenho ( Gráfico 4.2 ).
Nos dois quartos de menor desempenho, sem diferença estatisticamente significativa entre si, 36,2 \% no 1º quarto e 35,0 \% no 2º consideraram a prova difícil ou muito difícil.
Nos quartos de maior desempenho a proporção de alunos que julgaram a prova difícil ou muito difícil – 30,7 \% no 3º quarto e 24,4 \% no 4º quarto – foi menor do que nos demais e significativamente diferente entre si.
Para todos os quartos de desempenho a alternativa modal para esta pergunta foi médio, com 54,8 \%, 58,7 \%, 62,1 \% e 64,7 \% dos respondentes de \textbf{cada} um dos quartos, de 1 a 4 respectivamente, valores crescentes com o desempenho. 4.1.2 Componente de Conhecimento Específico Ao responderem à Questão 2 – “ Qual o grau de dificuldade desta prova na parte de Componente Específico? ” – 33,2 \% dos estudantes classificaram -na como difícil ou muito difícil.
Além disso, o Componente de Conhecimento Específico da prova foi considerado com grau de dificuldade médio por 59,1 \% dos alunos, ou seja, a maioria ( Gráfico 4.3, Gráfico 4.4 e, no Anexo II, a Tabela II.2 ).
A análise das respostas dos estudantes quanto ao grau de dificuldade do Componente de Conhecimento Específico da prova, agregados por Grande Região, mostra que a diferença entre a maior e a menor proporção de alunos que a avaliaram como difícil ou muito difícil é estatisticamente significativa : 37,3 \% na região Sul e 22,8 \% na Nordeste.
O percentual de alunos que classificaram o grau de dificuldade como médio, no Componente de Conhecimento Específico, variou de 56,9 \% a 66,1 \%, para as regiões Sul e Nordeste, respectivamente, novamente situação inversa dos que consideraram a prova como difícil ou muito difícil.
Considerando -se a avaliação da dificuldade das questões do Componente de Conhecimento Específico da prova, de acordo com o desempenho dos estudantes, observam -se diferenças estatisticamente significativas entre o quarto superior e os demais : a proporção dos que classificaram a parte específica como difícil ou muito difícil variou de 35,6 \% ( 1º quarto ) a 28,6 \% ( 4º quarto ), valores decrescentes com o desempenho.
Já a alternativa modal para a Questão 2 foi o grau médio, com 55,4 \% do quarto inferior e 62,4 \% do superior optando por esta resposta, valores crescentes com o desempenho. 4.2 EXTENSÃO DA PROVA EM RELAÇÃO AO TEMPO TOTAL Indagados quanto à extensão da prova, em relação ao tempo total oferecido para a sua resolução ( Questão 3 ), os estudantes apontaram, com maior incidência, a alternativas que considerava a extensão adequada, para todas as agregações consideradas ( Gráfico 4.5, Gráfico 4.6 e, no Anexo II, a Tabela II.3 ).
O percentual de alunos que responderam ser a extensão da prova adequada foi de 48,5 \%.
Já 43,3 \% dos inscritos presentes consideraram que o exame foi longo ou muito longo e 8,2 \% o avaliaram como curto ou muito curto.
Entre as Grandes Regiões a proporção daqueles que avaliaram a prova como longa ou muito longa em relação ao tempo total destinado à sua resolução variou de 38,2 \% na região Norte até 46,2 \% nas regiões Nordeste e Sudeste.
A proporção dos que consideraram a prova longa ou muito longa na região Norte é significativamente diferente da proporção dos que deram as mesmas respostas nas demais regiões, a exceção da região Centro-Oeste.
Considerando -se o desempenho dos alunos, nota -se ainda que o percentual dos que consideraram a extensão da prova adequada variou de 46,7 \%, no 4º quarto de desempenho, a 49,4 \% no 2º quarto de desempenho.
Observa -se, grosso modo, tendência decrescente conforme o desempenho.
No Gráfico 4.6, pode -se constatar que para os quatro quartos de desempenho a proporção de estudantes que consideraram a prova longa ou muito longa, em relação ao tempo total destinado à sua resolução, manteve -se em \textbf{torno} de 43 \%, ligeiramente crescente com o desempenho, mas sem diferenças estatisticamente significativas, com exceção do 4º quarto de desempenho com valor estatisticamente mais alto do que os demais. 4.3 COMPREENSÃO DOS ENUNCIADOS DAS QUESTÕES 4.3.1 Componente de Formação Geral Com relação aos enunciados das questões do Componente de Formação Geral ( Questão 4 ), as opiniões foram positivas, já que 68,7 \% dos alunos avaliados consideraram os enunciados de todas ou da maioria das questões claros e objetivos ( Gráfico 4.7, Gráfico 4.8 e, no Anexo II, a Tabela II.4 ).
Na análise regional, a percentagem de estudantes que avaliaram que todos ou a maioria dos enunciados das questões do Componente de Formação Geral \textbf{estavam} claros e objetivos variou de 67,1 \% na região Sudeste a 71,2 \% na região Nordeste, sendo esta diferença estatisticamente significativa.
A análise das percepções dos estudantes sobre a clareza e objetividade dos enunciados permite afirmar que todos ou a maioria dos enunciados de questões relativas ao Componente de Formação Geral foram considerados claros e objetivos pela a maior parte dos respondentes ( maior do que 67 \% em todas as regiões e maior do que 59 \% para todos os quartos de desempenho ).
Segundo o desempenho, observa -se que a proporção dos que emitiram esta opinião cresce conforme o desempenho aumenta, com diferenças estatisticamente significativas entre todos os quartos de desempenho.
No quarto superior, a clareza e objetividade de todos ou da maioria dos enunciados das questões foi percebida por 76,6 \% dos alunos e no quarto de desempenho inferior tal avaliação foi emitida por 59,7 \% deles. 4.3.2 Componente de Conhecimento Específico Com relação aos enunciados das questões do Componente de Conhecimento Específico da prova, para 68,4 \% dos estudantes avaliados da Área de Serviço Social a clareza e a objetividade ( Questão 5 ) \textbf{estavam} presentes em todas ou na maioria das questões ( Gráfico 4.9, Gráfico 4.10 e, no Anexo II, a Tabela II.5 ).
A maioria dos estudantes de todas as Grandes Regiões brasileiras considerou claros e objetivos todos ou a maioria dos enunciados das questões do Componente de Conhecimento Específico da prova, percentual sempre maior do que 67 \%.
As diferenças entre as regiões não são estatisticamente significativas, a não ser entre a região Nordeste ( com maior proporção ) e as regiões Sudeste, Sul e Centro-Oeste ( com as menores proporções ).
A proporção de estudantes que consideraram os enunciados das questões claros e objetivos apresenta uma tendência crescente em relação ao aumento de desempenho : mais elevada no quarto superior ( 77,3 \% ) se comparada ao quarto inferior de desempenho ( 59,6 \% ).
As diferenças dos resultados entre todos os quartos de desempenho são estatisticamente significativas. 4.4 SUFICIÊNCIA DAS INFORMAÇÕES/INSTRUÇÕES FORNECIDAS Ao avaliarem as informações/instruções fornecidas para a resolução das questões ( Questão 6 ), 74,0 \% dos respondentes da Área de Serviço Social de todo o Brasil afirmaram que estas eram até excessivas ou suficientes em todas ou na maioria das questões ( Gráfico 4.11, Gráfico 4.12 e, no Anexo II, a Tabela II.6 ).
Quanto à distribuição de respondentes pelas Grandes Regiões observa -se que a proporção de estudantes que consideraram as informações/instruções fornecidas até excessivas ou suficientes em todas ou na maioria das questões foi sempre superior a 73 \%, chegando a 77,3 \% na região Nordeste.
As diferenças entre as proporções das regiões não são estatisticamente significativas, a não ser entre a região Nordeste ( com a maior proporção ) e as regiões Sudeste, Sul e Centro-Oeste ( com as menores proporções ).
Levando -se em conta o desempenho dos participantes, notam -se diferenças estatisticamente significativas das opiniões de estudantes entre todos os quartos de desempenho, como mostra o Gráfico 4.12.
O percentual de participantes que avaliaram as informações/instruções como até excessivas ou suficientes em todas ou na maioria das questões foi mais elevado no quarto superior ( 82,3 \% ), percentual superior à média nacional ( 74,0 \% ).
Já no quarto inferior, a suficiência das informações/instruções declarada como até excessiva em todas ou na maioria das questões foi percebida por 65,5 \% dos respondentes, com um comportamento crescente com os quartos de desempenho. 4.5 DIFICULDADE ENCONTRADA AO RESPONDER À PROVA \textbf{Perguntados} sobre as dificuldades com as quais se depararam ao responder à prova ( Questão 7 ), 4,8 \% dos estudantes apontaram o desconhecimento do conteúdo.
Para 46,4 \%, a forma diferente de abordagem do conteúdo foi indicada como dificuldade.
Já a falta de motivação para fazer a prova foi a dificuldade apontada por 22,1 \% dos respondentes.
Considerando -se todo o Brasil, 19,7 \% dos respondentes afirmaram que não tiveram qualquer tipo de dificuldade para responder à prova ( Tabela II.7 no Anexo II ).
Os Gráficos 4.13 e 4.14 apresentam os percentuais de estudantes que apontaram o desconhecimento do conteúdo como dificuldade percebida ao responder à prova.
Na análise por Grandes Regiões, o percentual de inscritos e presentes que apontaram o desconhecimento do conteúdo como dificuldade ao responder à prova não superou 6 \%.
Os percentuais variaram de 3,1 \% na região Nordeste a 5,5 \% na Centro-Oeste, e esta diferença é estatisticamente significativa.
A forma diferente de abordagem do conteúdo foi a escolha modal dos estudantes, com percentuais que variaram de 34,5 \% ( região Nordeste ) a 51,6 \% ( região Sul ).
O percentual de alunos que citaram a falta de motivação como dificuldade variou de 16,6 \% ( região Norte ) a 27,5 \% ( região Nordeste ).
A proporção dos que declararam não ter qualquer dificuldade para responder à prova variou de 18,4 \% na região Sul a 25,1 \% na Nordeste.
Com relação aos quartos de desempenho, o desconhecimento do conteúdo, foi a opção escolhida por 3,8 \% dos estudantes do quarto superior e 6,8 \% do quarto inferior, valores decrescentes com o desempenho.
A alternativa modal para os alunos, quando agregados pelos quartos de desempenho, foi que a dificuldade encontrada foi causada pela forma diferente de abordagem do conteúdo : 46,8 \% no quarto inferior e 40,1 \% do quarto superior assim o responderam. 4.6 CONTEÚDOS DAS QUESTÕES OBJETIVAS DA PROVA Ao analisarem os conteúdos das questões objetivas da prova ( Questão 8 ), um percentual muito pequeno dos estudantes avaliados, apenas 5,4 \%, afirmou que não estudou ainda a maioria desses conteúdos ( Gráfico 4.15, Gráfico 4.16 e a Tabela II.8 no Anexo II ).
A maioria ( 65,2 \% ) afirmou ter estudado e aprendido muitos ou todos os conteúdos.
Na análise por Grande Região, a proporção de respondentes que escolheram a opção não estudou ainda a maioria desses conteúdos, foi pequena.
Observa -se que nas regiões Sul ( 6,3 \% ) e Centro-Oeste ( 6,0 \% ), apesar de pequenas, as proporções foram maiores do que a média nacional ( 5,4 \% ).
Dividindo as regiões em dois grupos, acima e abaixo da média nacional, em \textbf{cada} um dos dois grupos não existem diferenças estatisticamente significativas entre as regiões, no entanto na comparação entre esses grupos a diferença existe, quando não se considera a região Norte ( que não apresenta diferença estatisticamente significativa com as regiões Sul e Centro-Oeste ).
Em todas as regiões, a maioria dos presentes afirmou ter estudado e aprendido muitos ou todos os conteúdos, com proporções variando entre 58,9 \% na região Centro-Oeste e 75,5 \% na Nordeste.
Considerando -se separadamente as opiniões de estudantes dos quatro quartos de desempenho, observa -se que, no quarto inferior, 10,1 \% ofereceram como resposta que não estudou ainda a maioria desses conteúdos, sendo 1,8 \% os do quarto superior com a mesma resposta, com valores decrescentes com o desempenho.
A diferença entre os alunos que optaram por este motivo de dificuldade é estatisticamente significativa entre todos os quartos de desempenho.
Tendo em conta o quarto superior, 81,1 \% dos alunos afirmaram ter estudado e aprendido muitos ou todos os conteúdos.
No outro extremo, no quarto inferior, 51,4 \% dos alunos responderam da mesma forma, com valores crescentes com o desempenho. 4.7 TEMPO GASTO PARA CONCLUIR A PROVA Ao responderem sobre o tempo de conclusão da prova ( Questão 9 ), um pouco mais de 2/3 dos estudantes ( 68,7 \% ) afirmou ter gasto entre duas e quatro horas ( Gráfico 4.17, Gráfico 4.18 e, no Anexo II, a Tabela II.9 ).
Considerando -se as cinco Grandes Regiões brasileiras, nas regiões Norte ( 68,4 \% ), Nordeste ( 66,0 \% ) e Sudeste ( 65,5 \% ) o percentual dos que utilizaram entre duas e quatro horas para finalizar a prova foi inferior à média nacional ( 68,7 \% ).
Em duas das cinco Grandes Regiões, Sul e Centro-Oeste, o percentual de alunos que dispensaram entre duas e quatro horas para concluir a prova, superou 70 \%, como mostra o Gráfico 4.17.
Uma vez sendo analisadas as alternativas escolhidas pelos estudantes que se situam nos diferentes quartos de desempenho, observa -se que uma proporção bem maior de participantes no quarto superior declarou ter gasto entre duas e quatro horas para concluir a prova quando comparados com os do quarto inferior, respectivamente 73,1 \% e 61,9 \%.
As diferenças entre os quatro quartos são estatisticamente significativas, exceto entre o 2º e o 3º quartos, e evidenciam uma tendência crescente com o desempenho.

% matched lemmas: cada, estar, lembrar, perguntar, torno
\end{textsample}
